\section{Proposed Methodology: DS-HCRS Framework}

\textbf{PLANNER}

(Conceptual - no implementation details.)

\subsection{Design Philosophy}

Principles:

\begin{itemize}
    \item human-aware optimization
    \item adaptive cognition modeling
    \item long-term value
\end{itemize}
    
\subsection{System Overview Architecture}

High-level diagram description:

User → State Modeling → Dual Systems → Gating → Recommendation → Feedback.

\subsection{System 1: Fast Engagement Model}

Role:

\begin{itemize}
    \item reactive ranking
    \item immediate reward modeling.
\end{itemize}
    
\subsection{Fatigue \& Cognitive Burden Model}

Define variables:

\begin{itemize}
    \item repetition
    \item entropy decay
    \item session length
    \item novelty loss
\end{itemize}    

Explain theoretical grounding.

\subsection{System 2: Reflective Long-Term Optimizer}

Conceptual RL formulation.

Explain:

\begin{itemize}
    \item delayed reward
    \item retention proxy
    \item welfare regularization.
\end{itemize}    

\subsection{Dual-System Gating Mechanism}

Decision policy rationale.

Explain adaptive reasoning depth idea.

\subsection{Optional Concept Embedding Module}

LLM embeddings for cold start.

\subsection{Learning Objective Formulation}

General reward:

\[
R_t = \alpha E_t + \beta L_t - \gamma F_t
\]

Explain meaning conceptually.

\subsection{Expected Advantages Over Existing Systems}

Table comparison:

\begin{table}[!t]
\caption{An Example of a Table\label{tab:table1}}
\centering
\begin{tabular}{|l||c||c|}
\hline
Feature & Traditional & Proposed \\
\hline
Fatigue modeling & cross & tick \\
Long-term reward & Partial & Explicit \\
Adaptive reasoning & cross & tick \\
\hline
\end{tabular}
\end{table}